% BHU PYQ -- Transcribed & Corrected
% Paper: SANMJ-32 / SANMJ32: Vyakaran-II (Vyakaran Evam Bhasha Vigyan)
% Exam: B.A. (Hons.) Semester III Examination, 2025-26 (NEP)
% Category: Major Core
% Department: Department of Sanskrit, Faculty of Arts, Banaras Hindu University

\documentclass[11pt,a4paper]{article}
\usepackage[a4paper,top=2.5cm,bottom=2.5cm,left=2.8cm,right=2.8cm,headheight=14pt]{geometry}
\usepackage{amsmath,amssymb}
\usepackage[T1]{fontenc}
\usepackage[utf8]{inputenc}
\usepackage{lmodern,microtype,fancyhdr,enumitem,hyperref,lastpage,parskip}

\hypersetup{
    colorlinks=true,
    urlcolor=black,
    linkcolor=black,
    pdftitle={SANMJ32: Vyakaran II (Vyakaran Evam Bhasha Vigyan) -- BHU B.A. Sem III 2025-26 (NEP)}
}

\pagestyle{fancy}
\fancyhf{}
\renewcommand{\headrulewidth}{0.4pt}
\renewcommand{\footrulewidth}{0.4pt}
\fancyhead[L]{\small \textit{Banaras Hindu University}}
\fancyhead[R]{\small \textit{SANMJ-32: Vyakaran II (Sem III 2025-26)}}
\fancyfoot[C]{\small Page \thepage\ of \pageref{LastPage}}

\newlist{parts}{enumerate}{2}
\setlist[parts,1]{label=(\alph*),leftmargin=*,itemsep=4pt,topsep=2pt}

\newcommand{\pts}[1]{\hfill \textit{[#1]}}

\begin{document}

\begin{center}
  {\large\bfseries BANARAS HINDU UNIVERSITY}\\[2pt]
  {\normalsize B.A. (Hons.) Semester III Examination, 2025-26 (NEP)}\\[6pt]
  \rule{0.8\linewidth}{0.5pt}\\[6pt]
  {\large\bfseries DEPARTMENT OF SANSKRIT}\\[4pt]
  {\large\bfseries Paper Code: SANMJ-32 : व्याकरण-II (व्याकरण एवं भाषा विज्ञान)}\\[2pt]
  {\small Ref. Code: 2025-26/D-III/056}\\[4pt]
  {\normalsize Time: 2:30 Hours \hfill Full Marks: 70}
\end{center}

\rule{\linewidth}{0.4pt}

\smallskip

\noindent \textit{(Write your Roll No. at the top immediately on the receipt of this question paper)}\\[2pt]
\noindent \textit{सूचना : सभी प्रश्नों के उत्तर निर्देशानुसार दीजिए।}

\bigskip

\begin{center}
  {\large\bfseries SECTION A / खण्ड-अ}\\[2pt]
  {\normalsize (Long Answer Type Questions / दीर्घ उत्तरीय प्रश्न)}
\end{center}

\noindent \textit{सूचना : किन्हीं तीन प्रश्नों के उत्तर दीजिए। प्रत्येक प्रश्न 10 अंकों का है।}\\[2pt]

\bigskip

\noindent \textbf{Question 1.} निम्नलिखित में से किन्हीं तीन सूत्रों की सोदाहरण व्याख्या कीजिए:\pts{10}
\begin{parts}
  \item प्रातिपदिकम्‌
  \item आदेच उपदेशेऽशिति
  \item सर्वादीनि सर्वनामानि
  \item एकाचो बशो भष्‌ झषन्तस्य स्ध्वो:
  \item नलोपः प्रातिपदिकान्तस्य
\end{parts}

\medskip\rule{\linewidth}{0.2pt}\medskip

\noindent \textbf{Question 2.} कारक प्रकरण के आधार पर 'अकथितं च' सूत्र की सोदाहरण व्याख्या कीजिए। \pts{10}

\medskip\rule{\linewidth}{0.2pt}\medskip

\noindent \textbf{Question 3.} भाषा विज्ञान के अनुसार ध्वनियों के वर्गीकरण के प्रमुख आधारों की विवेचना कीजिए। \pts{10}

\medskip\rule{\linewidth}{0.2pt}\medskip

\noindent \textbf{Question 4.} भारोपीय भाषा परिवार की प्रमुख विशेषताओं एवं शाखाओं का सचित्र/सोदाहरण वर्णन कीजिए। \pts{10}

\medskip\rule{\linewidth}{0.2pt}\medskip

\noindent \textbf{Question 5.} ध्वनि परिवर्तन के प्रमुख कारणों पर विस्तृत प्रकाश डालिए। \pts{10}

\bigskip
\rule{\linewidth}{0.2pt}
\bigskip

\begin{center}
  {\large\bfseries SECTION B / खण्ड-ब}\\[2pt]
  {\normalsize (Short Answer Type Questions / लघु उत्तरीय प्रश्न)}
\end{center}

\noindent \textit{सूचना : किन्हीं चार प्रश्नों के उत्तर दीजिए। प्रत्येक प्रश्न 5 अंकों का है।}\\[2pt]

\bigskip

\noindent \textbf{Question 6.} 'कर्मणि द्वितीया' सूत्र की सोदाहरण व्याख्या कीजिए। \pts{5}

\medskip\rule{\linewidth}{0.2pt}\medskip

\noindent \textbf{Question 7.} 'स्वतन्त्रः कर्ता' सूत्र को स्पष्ट कीजिए। \pts{5}

\medskip\rule{\linewidth}{0.2pt}\medskip

\noindent \textbf{Question 8.} ग्रिम नियम (Grimm's Law) की सोदाहरण व्याख्या कीजिए। \pts{5}

\medskip\rule{\linewidth}{0.2pt}\medskip

\noindent \textbf{Question 9.} वर्नर नियम (Verner's Law) का संक्षिप्त विवरण दीजिए। \pts{5}

\medskip\rule{\linewidth}{0.2pt}\medskip

\noindent \textbf{Question 10.} अर्थ परिवर्तन की प्रमुख दिशाओं की सोदाहरण विवेचना कीजिए। \pts{5}

\medskip\rule{\linewidth}{0.2pt}\medskip

\noindent \textbf{Question 11.} संस्कृत भाषा की प्रमुख भाषावैज्ञानिक विशेषताओं को रेखांकित कीजिए। \pts{5}

\bigskip
\rule{\linewidth}{0.2pt}
\bigskip

\begin{center}
  {\large\bfseries SECTION C / खण्ड-स}\\[2pt]
  {\normalsize (Very Short Answer Type Questions / अतिलघु उत्तरीय प्रश्न)}
\end{center}

\noindent \textit{सूचना : निम्नलिखित सभी दस प्रश्नों का उत्तर एक वाक्य में दीजिए। प्रत्येक प्रश्न 2 अंकों का है।}\\[2pt]

\bigskip

\noindent \textbf{Question 12.}
\begin{parts}
  \item 'सर्वा' इस पद में कौन-सी विभक्ति एवं वचन है? \pts{2}
  \item 'मह्यम्' शब्द रूप किस सर्वनाम का है और किस विभक्ति का रूप है? \pts{2}
  \item 'राम' शब्द का पञ्चमी विभक्ति एकवचन में क्या रूप बनता है? \pts{2}
  \item 'अकथितं च' सूत्र के द्वारा किस संज्ञा का विधान होता है? \pts{2}
  \item 'हरये रोचते भक्ति:' में 'हरये' पद में किस सूत्र से चतुर्थी विभक्ति हुई है? \pts{2}
  \item भाषा विज्ञान के जनक के रूप में किसे जाना जाता है? \pts{2}
  \item भारोपीय भाषा परिवार की शतम एवं केन्तुम वर्गों में विभाजन का आधार क्या है? \pts{2}
  \item 'ग्रामं गच्छन् तृणं स्पृशति' में 'तृणम्' की कर्म संज्ञा किस सूत्र से हुई है? \pts{2}
  \item अष्टाध्यायी में कुल कितने अध्याय और पाद हैं? \pts{2}
  \item 'सद्विद्या' शब्द में सन्धि-विधान करने वाला सूत्र लिखिए। \pts{2}
\end{parts}

\vfill
\rule{\linewidth}{0.4pt}

\begin{center}\small
\href{https://bhu-exam.vercel.app/}{Click here to visit our website}\\[4pt]
\href{https://me-aryan.vercel.app/}{Click here to visit our contributor}
\end{center}

\end{document}
